\scanpage{5}
d'une action de $\mathfrak{S}_3$ sur \inserted{le $\Lambda$-module} $\Lambda^2$ qu'on peut expliciter ainsi. On considère l'action de $\mathfrak{S}_3$ sur
\[
\Lambda^3/\underbrace{\Lambda}_{\text{diagonale}}=V'_{\{1,2,3\}}(\Lambda)={}_{2}V'_{\Lambda}\tag*{(6)}
\]
de sorte que l'on a un isom canonique de $\mathfrak{S}_3$-ensembles
\[
\overset{\unsure{\sim|\deleted{\sim}|-}}{B}(M)\simeq\mathrm{Isom}_{\Lambda}(V'_{\Lambda},M)\tag*{(7)}
\]
D'ailleurs \deleted{\unsure{l'is|l'i}} l'application (5) est composée de (7) et de la bijection
\[
\mathrm{Isom}_{\Lambda}(V'_{\Lambda},M)\xrightarrow{\sim}\mathrm{Isom}_{\Lambda}(\Lambda^2,M)
\]
déduit de l'isomorphisme canonique
\[
\left\{\begin{aligned}
\Lambda^2&\xrightarrow{\sim}\Lambda^3/\Lambda\\
(x,y)&\longmapsto(x,y,0)\bmod\Lambda
\end{aligned}\right. \ .\tag*{(8)}
\]
Revenons à un $(x,y,z)\in B(M)$, notons qu'on \unsure{aura|a} alors
\[
x\wedge y=y\wedge z=z\wedge x\ \in
\underbrace{\det(M)^*}_{\substack{\text{ens des}\\\text{bases du }\Lambda\text{-module}\\\text{inversible }\det M=\Lambda^2 M}}\tag*{(9)}
\]
De façon intrinsèque, si $A\in\widetilde F(M)$, \deleted{choisissant un ordre total $(x,y,z)$} \inserted{\unsure{l'écrivant|écrivant} $A=\{x,y,z\}$},
% archive sheet number: 4
